% ================================================================
\section{The Epistemology of Machine-Verified Proof}
\label{sec:epistemology}
% ================================================================

\subsection{Three Conceptions of Proof}

The history of the philosophy of mathematics offers three
progressively refined conceptions of what it means to prove
a mathematical claim.

\subsubsection{Proof as Persuasion}

In the classical tradition from Euclid through the 19th century,
a proof is a rhetorical act: a chain of reasoning designed to
compel assent from a competent reader. Hardy~\cite{hardy1940}
described the mathematician's task as presenting a ``proof of the
theorem'' that makes the theorem ``seem almost obvious'' to the
trained eye. Thurston~\cite{thurston1994} refined this into an
explicitly social account: mathematical proof is ``a social process
by which mathematicians come to agree.''

The strengths of this conception are well known. The weaknesses
became apparent in the 20th century, as proofs grew too long
(the classification of finite simple groups exceeds 10{,}000 pages)
or too computationally intensive (the four-color theorem,
Hales's proof of the Kepler conjecture) for any single human mind
to verify. Lakatos~\cite{lakatos1976} showed that even
``elementary'' proofs can contain unnoticed gaps that survive
decades of expert scrutiny.

\subsubsection{Proof as Formal Derivation}

Hilbert's program proposed that mathematical proof should be
reduced to mechanical symbol manipulation: a proof is a finite
sequence of formulas, each either an axiom or derived from earlier
formulas by specified rules, whose last formula is the theorem.
This conception eliminates the need for human judgment---in
principle, a clerk with no mathematical training could verify
the proof by checking each rule application.

The Curry--Howard correspondence deepened this to an identification:
proofs ARE programs, and the propositions they prove are the types
of those programs. In this view, a proof of $\forall x,\, P(x)$
is a function that, given any $x$, produces evidence for $P(x)$.
Verification becomes type-checking: does this program (proof) have
the claimed type (theorem)?

Modern proof assistants---Lean, Coq, Isabelle, Agda---implement
this identification literally. A ``proof'' in Lean~4 is a term in
the Calculus of Inductive Constructions (CIC) that type-checks
against a specification. The verification is performed by a small,
trusted kernel (approximately 10{,}000 lines of C++ for Lean~4)
whose correctness is independently auditable.

\subsubsection{The Cathedral's Synthesis}

The Cathedral occupies an unusual position in this landscape.
It is a formal proof in the Hilbert--Curry--Howard sense: it
type-checks, and the Lean kernel certifies its correctness. But
what it proves is not the Riemann Hypothesis. It proves:

\begin{enumerate}
  \item \textbf{The converse} ($d_N^2 \to 0 \implies \RH$):
    with zero custom axioms.
  \item \textbf{The forward direction} ($\RH \implies d_N^2 \to 0$):
    conditional on one literature axiom.
  \item \textbf{The equivalence} ($\RH \iff d_N^2 \to 0$):
    the conjunction of (1) and (2), conditional on the same axiom.
\end{enumerate}

This creates a \emph{gradient of certainty}:
\begin{center}
\small
\begin{tabular}{@{}lll@{}}
\toprule
\textbf{Claim} & \textbf{Axioms} & \textbf{Epistemic Status} \\
\midrule
$d_N^2 \to 0 \implies \RH$ & 0 & Machine-certified \\
$\RH \implies d_N^2 \to 0$ & 1 (literature) & Conditional \\
$\RH \iff d_N^2 \to 0$ & 1 (literature) & Conditional \\
RH itself & = Crown Axiom & Open \\
\bottomrule
\end{tabular}
\end{center}

\noindent What kind of knowledge is this?

\subsection{Conditional Knowledge and Its Value}

Avigad~\cite{avigad2006} has argued that the value of a
mathematical proof lies not merely in establishing truth, but in
providing \emph{understanding}---a grasp of the logical
relationships between concepts. By this criterion, the Cathedral
provides genuine mathematical knowledge: it establishes, with
machine-checked certainty, the precise logical relationship between
the Riemann Hypothesis and a concrete analytic inequality.

Consider an analogy from physics. Newton's law of gravitation
does not \emph{prove} that the planets orbit the sun. It
establishes a conditional: \emph{if} gravity obeys an inverse-square
law, \emph{then} the planetary orbits are conic sections. The
conditional itself is a theorem; the hypothesis (inverse-square
gravity) is an empirical claim, verified to high precision but
never ``proved'' in the mathematical sense. No physicist considers
Newton's work incomplete because the inverse-square law remains,
strictly speaking, an axiom.

The Cathedral stands in the same relationship to the Riemann
Hypothesis. The conditional equivalence $\RH \iff d_N^2 \to 0$ is
a theorem---fully proved, machine-verified, not dependent on RH
being true. The Crown Axiom (the Gram form bound) is the empirical
hypothesis, verified computationally to $N = 55{,}440$ with
31-digit precision. Just as Newton's conditional is valuable
regardless of whether gravity is ``truly'' inverse-square, the
Cathedral's conditional is valuable regardless of whether RH
is ``truly'' true.

\subsection{The Gradient of Certainty}

The Cathedral exhibits a remarkable epistemological structure: a
\emph{gradient} from absolute certainty to complete uncertainty,
traversed in precisely identified steps.

\begin{enumerate}
  \item \textbf{Lean kernel axioms} (propext, Classical.choice,
    Quot.sound): accepted by the entire Lean community. Certainty
    level: foundational.

  \item \textbf{Mathlib infrastructure} (Stirling, harmonic numbers,
    Cauchy integral formula): proved from kernel axioms, machine-checked.
    Certainty level: absolute (modulo kernel trust).

  \item \textbf{Graduated theorems} (PNT Möbius sums, Abel engine,
    NB distance formula, bridges): formerly axioms, now proved.
    Certainty level: absolute.

  \item \textbf{PNT axioms} (mu\_pnt\_alt, pnt\_mu\_log\_sq\_div\_k):
    provable from the PrimeNumberTheoremAnd project, which has
    formalized these results independently. Certainty level:
    very high (will become absolute when imported).

  \item \textbf{The Crown Axiom} ($v^TGv \leq 1 + K/\ln N$): this
    IS the Riemann Hypothesis, reformulated. Certainty level:
    the central open question of mathematics.
\end{enumerate}

The graduated gates---the axioms that were successively converted
to theorems (56 $\to$ 42 $\to$ 7 $\to$ 4 $\to$ 2 $\to$ 1)---are
the most philosophically significant feature of this gradient. Each
graduation was an act of converting faith into knowledge: a claim
that was previously accepted on the authority of classical
mathematics was replaced by a machine-verified derivation. The
progression is not merely a bookkeeping achievement; it is an
\emph{epistemological compression}, concentrating the entire
unknown content of the Riemann Hypothesis into a single, precisely
stated inequality.

\subsection{The False Axiom and the Limits of Intuition}

A particularly instructive episode occurred during the axiom
graduation process. The axiom \lean{remainder\_bound\_abstract},
which asserted an operator norm bound on the off-diagonal correction
matrix $E_{\mathrm{ratio}}$, was believed to be true on the basis
of mathematical intuition and partial analysis. Numerical
experiments (conducted by one of the AI collaborators) then
demonstrated that the axiom was \textbf{false} for general
vectors---the operator norm grows logarithmically, not decaying
as claimed.

This discovery carries philosophical weight. It demonstrates that:
\begin{enumerate}
  \item Mathematical intuition, even when grounded in structural
    analysis, can be wrong.
  \item Machine computation can detect errors that human reasoning
    misses.
  \item The Lean compiler, which treats axioms as opaque declarations,
    offers no protection against false axioms---it will cheerfully
    derive nonsense from false assumptions.
\end{enumerate}

The resolution was equally instructive. The false global bound was
replaced by three \emph{entry-level} theorems about the individual
matrix entries: each $E_{\mathrm{ratio}}(j,k)$ is non-positive
(proved from the structural identity $(a-b)\ln(b/a) \leq 0$), and
each has magnitude bounded by $(b-a)^2/(2a^2b)$. These entry-level
bounds are both \emph{stronger} and \emph{true}---they provide
pointwise control that the false global bound could not.

The philosophical lesson is that formalization disciplines
mathematical intuition in a way that informal reasoning cannot.
The false axiom survived months of informal analysis; it fell in
minutes to numerical experiment. The replacement theorems, once
discovered, were proved in Lean with zero sorry. The episode
vindicates the formalization project's methodology: state your
assumptions explicitly, test them computationally, prove them
formally, and be prepared to discover that some of them are wrong.

\subsection{The Tri-Crown and Epistemological Pluralism}

The Cathedral offers not one but three independent paths to the
Riemann Hypothesis:
\begin{itemize}
  \item The \textbf{Analytic Crown}: RH via one literature axiom
    (Baez-Duarte's forward direction).
  \item The \textbf{Cybernetic Crown} (Oracle Bridge): RH via one
    computation axiom (GPU-certified Gram matrices at highly
    composite numbers).
  \item The \textbf{Gram Crown}: RH via one arithmetic inequality
    (the Gram form bound along a subsequence).
\end{itemize}

All three paths share the same zero-axiom converse. They differ
only in how they establish the forward direction. This architectural
feature has epistemological significance: it demonstrates that the
truth of RH, if it is true, can be \emph{witnessed} from multiple
independent perspectives, each with its own epistemic character.

The Analytic Crown relies on the authority of a published journal
article. The Cybernetic Crown relies on the correctness of a
computational pipeline. The Gram Crown relies on the truth of a
pure arithmetic inequality. These are three fundamentally different
\emph{kinds} of evidence. The Cathedral's architecture makes their
equivalence precise: all three are logically sufficient, and the
machine-verified infrastructure converts each into a proof of RH.

This is a form of mathematical \emph{perspectivism}: there is no
single ``correct'' proof of RH, only multiple valid perspectives
on the same underlying truth. The Cathedral makes this pluralism
explicit and machine-verifiable.
